\documentclass[a4paper,10pt]{amsbook}
\usepackage{fullpage}
\usepackage{amsmath,amsfonts,amssymb}
\usepackage{graphicx}
\usepackage{newverbs}
\usepackage{fancyvrb}
\usepackage{hyperref}
\usepackage{tikz}
\usepackage{enumitem}

\input{macros}

% Create a new list environment
\newlist{secenum}{enumerate}{1}

% Configure it to use the custom counter without resetting it at \begin{secenum}
\setlist[secenum,1]{
    label=\thesection.\arabic*,
    ref=\thesection.\arabic*
}


\begin{document}
\begin{center}
 {\Large Maths with Python --- Python reference}
\end{center}
\vspace{4ex}

\section{Getting started}
\label{sec-start}

These instructions will assume that you are comfortable with using
Python in Jupyter notebooks in VS Code, and that you have installed 
the packages \pcode+sympy+, \pcode+numpy+, \pcode+scipy+ and
\pcode+plotly+ using \pcode+pip+ or \pcode+conda+.  You will also 
need to import these packages and set some options.  The most
convenient way to do this is as follows.  Distributed together with
this file is a directory called \pcode+header+ containing a single
file \pcode+__init__.py+.  You should copy this directory to be a
subdirectory of the directory where you keep your notebooks for this
course, then you should include the line 
\begin{pcodeblock}
 from header import *
\end{pcodeblock}
as the top line of each of your notebooks.

\section{Notation for variables}

\begin{secenum}
 \item\label{symbols} By default names like $x$ refer to variables as used
  in programming languages like Python.  They are expected to have a
  value (which might be a number or a list or something more
  complicated).  If you try to use a variable (for example, by trying
  to expand $(1+x)^4$) without assigning a value  then you will get an
  error.  If you instead want to treat $x$ as an abstract symbol 
  then you need to enter \pcode+x = sp.symbols('x')+.  You can do this for 
  several variables at the same time using syntax like 
  \pcode+x,y,a,b = sp.symbols('x y a b')+.  
%
 \item\label{subscripts} If you want variables 
  with indices like $m_0,m_1,m_2$ then you should enter 
  \pcode+m = sp.IndexedBase('m')+, and then enter \pcode+m[2]+ for
  $m_2$, for example.
%
 \item\label{variables} Most variables will have single-letter
  names, such as $x$ or $A$. 
%
  \item\label{case} Note that case is significant: Python regards
  $A$ as different from $a$, and $x$ as different from $X$. 
%
  \item\label{longvars} It is perfectly permissible to use long
  names for variables, such as \pcode~eqns~ or \pcode~sols~ or
  \pcode~SpeedOfLight~. 
%
  \item\label{reserved} However, you may not use words that
  already have a special meaning in Python.  In particular,
  you cannot call a variable \pcode~lambda~. 
%
  \item\label{greek} If you enter \pcode+mu = sp.symbols('mu')+ then in some
  (but not all) contexts, \pcode+mu+ will be displayed as $\mu$.
%  Alternatively, there are various ways to set things up so that you
%  can just type the symbol \pcode+μ+ (or cut and paste it from
%  somewhere).  Then you can enter \pcode+μ = sp.symbols('μ')+ and 
%  \pcode+μ+ will always be displayed as $\mu$.
\end{secenum}

\section{Generalities}

\begin{secenum}
 \item\label{help} If you know the name of a function that you want
 to use (for example, \pcode+sp.solve+) you can enter
 \pcode+help(sp.solve)+ to get some documentation.  You can also get
 a shorter version of that documentation by just typing
 \pcode+sp.solve+ in a code cell and hovering over it with your mouse.
%
 \item\label{problems} See Section~\ref{sec-problems}
 for advice on some common problems. 
%
 \item\label{lastoutput} If you have several commands in the same code cell
 then Python will only print the result of the last one by default.
 If you want to see other things then you should add explicit
 \pcode+print()+ or \pcode+display()+ commands.
%
 \item\label{semicolon}
 If you end a command with a semicolon then Python will not 
 print the result even if it is the last command in the cell.
 You can also achieve the same effect by entering \pcode+None+
 as the last command in the cell.
%
 \item\label{underscore} You can use the symbol \pcode~_~
 to refer to the last thing that Python calculated and displayed,
 and \pcode~__~ to refer to the one before that, and so on. 
 However, this mechanism cannot be used to refer to things that were
 calculated but not displayed.
%
 \item\label{lasteval}
 If you move around a worksheet inserting things in
 different places, you should remember that \pcode~_~ refers
 to the \textbf{most recent} result, which may not be the
 one that you see directly above the line where you are
 typing.  
%
 \item\label{equals}
 It is important to understand the difference between
 \pcode~=~ and \pcode~==~ and \pcode+sp.Eq()+.  
 \begin{itemize}
  \item If we enter $x=3$, this defines $x$ to be equal to $3$, so
   Python will replace $x$ by $3$ whenever it sees an $x$.
  \item If we instead enter $x==3$, we are checking whether $x$ is
   currently straightforwardly equal to $3$.  If $x$ is an abstract
   symbol which has not been assigned a numerical value then $x==3$
   will just be false, even if $x$ could be equal to $3$ in some
   contexts.  
  \item We can instead enter \pcode+sp.Eq(x=3)+ to represent $x=3$
   as an abstract equation which might be true or false in different
   circumstances.  This is the right thing to do if we are setting up
   some equations that we later want to solve.
 \end{itemize}
%
 \item\label{funcdef}
 If you want to define a function (say $f(x)=x^2+3$) it does
 \emph{not} work to enter \pcode~f(x)=x^2+3;~.  See
 Section~\ref{sec-funcdef}. 
%
 \item\label{restart}
  Near the top of your VS Code window you should see a menu item
  marked \pcode+Restart+.  If you click this then Python will restart
  and all values that you might have assigned to variables will be
  removed.  You might want to do this at the beginning of each new
  exercise, otherwise stray values left over from previous exercises
  can cause trouble and confusion.  Note, however, that most notebooks
  will have a cell at the top with some \pcode+import+ statements.
  Restarting removes all imports, so you will need to go back to your
  import cell and execute it again before continuing.
%
 \item\label{unassign}
  If you just want to clear a single variable (say $x$) you
  can enter \pcode~del x~ instead.  To clear several variables, you
  can enter something like \pcode+del A,x,U+.
\end{secenum}

\section{Numerical approximation}

\begin{secenum}
 \item\label{float} To get an explicit numerical approximation for an
  expression, use the function \pcode~sp.N()~ function.  For example,
  \pcode~sp.N(sp.pi)~ gives $3.14159265358979$, and
  \pcode~sp.N(x/3)~ gives $0.3333333333 x$. 
%
 \item\label{floatdigits} The function \pcode~sp.N()~ generally gives an
  answer to $15$ significant figures.  If you want to
  calculate $\pi^2/6$ to $50$ digits (for example), you can
  enter \pcode~sp.N(sp.pi**2/6, 50)~.
%
 \item\label{digitslost} It is possible to lose accuracy for non-obvious
 reasons.  For example, if you enter \pcode+sp.N(1/7, 100)+ then
 Python will work from the inside out, first calculating $1/7$ to low
 accuracy and only then calling \pcode+sp.N()+ with a request for high
 accuracy that can no longer be achieved.  One way to fix this is to
 enter \pcode+sp.N(1,100)/7+, which requests high accuracy before
 attempting to divide by $7$.
\end{secenum}


\section{Algebraic operations}

\begin{secenum}
 \item\label{star} Use a \pcode~*~ for multiplication; for example,
 $2xy$ should be entered as \pcode~2*x*y~, and $(2x+1)\sin(3x)$
 should be entered as \pcode~(2*x+1)*sp.sin(3*x)~. 
%
 \item\label{twoletters}
 Note also that Python will treat \pcode~xy~ as a two-letter
 name for a single variable; if you want $x$ times $y$, you
 must enter \pcode~x*y~. 
%
 \item\label{timesbrackets}
 Just as in standard mathematical notation, bracketing is
 important.  To multiply $x+2$ by $x-2$, you must enter
 \pcode~(x+2)*(x-2)~, not \pcode~x+2*x-2~.  
%
 \item\label{powers} Enter powers using \pcode~**~,
 for example \pcode~x ** 6~ for $x^6$.  The notation \pcode+x^6+ does
 have a meaning (the exclusive or operator applied to $x$ and $6$) but
 it is almost never what you want; this can unfortunately lead to some
 confusing error messages.
%
 \item\label{powerbrackets}
 Just as in standard mathematical notation, bracketing is
 important.  To raise $x+2$ to the power $n+5$, you must
 enter \pcode~(x+2)**(n+5)~, not \pcode~x+2**n+5~.  To get the
 cube root of $x$, enter \pcode~x**(1/3)~, not
 \pcode~x**1/3~.  
%
 \item\label{sqrt} The square root of $x$ can be entered as
 \pcode~sp.sqrt(x)~.  Similarly, $\sqrt{b^2-4ac}$ is
 \pcode~sp.sqrt(b**2-4*a*c)~.  For more general roots you might be
 tempted to enter something like \pcode~(x+y)**(1/3)~ for the cube
 root of $x+y$.  However, it is actually necessary to enter 
 \pcode~(x+y)**sp.Rational(1,3)~ instead, to ensure that we use the
 exact fraction $1/3$ instead of converting it to an approximate
 decimal, which is what Python would do by default.
\end{secenum}

\section{Algebraic manipulation}

Suppose we have a complicated expression called $A$, which we want to
simplify or otherwise manipulate. 
\begin{secenum}
 \item\label{subs} To substitute for some or all of the variables in
  $A$, use the \pcode~subs~ command.  For example, if $A$ is
  $(a+b)^n$, you can change the $b$ to $1$ and the $n$ to
  $p/q$ by entering \pcode~A.subs({b:1,n:p/q})~, which gives
  $(a+1)^{p/q}$.  (Note that the correct syntax here is \pcode~b:1~,
  not \pcode~b=1~.)
%
 \item\label{solvesubs}
  Often the equations to be used in \pcode~subs(...)~ will
  come from the \pcode~solve(...)~ command.  For example, to
  find the value of $x^2+y^2$ at the point where $x+y=6$ and
  $x-y=2$, enter \pcode~sol=sp.solve([sp.Eq(x+y,6),sp.Eq(x-y,2)])~ 
  and then \pcode~(x**2 + y**2).subs(sol)~.
%
 \item\label{expand} Enter \pcode~sp.expand(A)~ to expand everything out.  For
  example, enter \pcode~sp.expand((x+y)**3)~ to expand
  $(x+y)^3$, giving $x^3+3x^2y+3xy^2+y^3$. 
%
 \item\label{simplify} To simplify $A$, try entering
  \pcode~sp.simplify(A)~.  If this 
  does not do what you want, you can try \pcode~sp.factor(A)~. 
  If $A$ involves trigonometric functions, try
  \pcode~sp.trigsimp(A)~ or \pcode~sp.expand_trig(A)~. 
%
 \item\label{symbolic} Note that rules like $(a^4)^{1/4}=a$ are
  \textbf{not valid} when $a$ is negative.  The function
  \pcode~sp.simplify()~ will not simplify expressions like this by
  default.  However, if you declare $x$ to be positive using
  \pcode~x = sp.symbols('x', positive=True)~ (instead of just 
  \pcode~x = sp.symbols('x')~) then sympy will apply rules that are
  only valid for positive $x$.  For example,
  \pcode+(x ** 4) ** sp.Rational(1,4)+ will simplify to $x$.  (However, 
  \pcode+(x ** 4) ** (1/4)+ will only simplify to $x^{1.0}$, which will 
  not simplify to $x$, because the exponent is only approximately equal 
  to one and not exactly equal.)  There are also more complex ways to
  specify that $x$ is an integer or lies in a certain range or has
  other useful properties that could be used when simplifying; see the
  sympy documentation about assumptions.
%
 \item\label{testequal} If you suspect that $A$ is equal to some other expression (say
  $B$) and want to check this, the best way is to enter
  \pcode~sp.simplify(A-B)~ and see if you get zero. 
%
 \item\label{collect} To collect together all the terms in $A$ involving the same
  power of $x$, enter \pcode~A.collect(x)~.  For example, this converts
  $1+ax+bx+cx^2+dx^2$ to $1+(a+b)x+(c+d)x^2$. 
%
 \item\label{coeff}
  If you just want the
  coefficient of $x^2$ in $A$ then you can instead enter \pcode~A.coeff(x,2)~
  or \pcode~A.coeff(x**2)~.  To get the constant term, you
  must enter \pcode~A.coeff(x,0)~ --- the syntax
  \pcode~A.coeff(x**0)~ will not work. 
\end{secenum}

\section{Constants}

\begin{secenum}
 \item\label{rational} If you enter a fraction like \pcode~2/9~ then it
  will immediately be converted to an approximate decimal, which is 
  not usually what you want.  If you want to treat it as an exact 
  fraction you should enter \pcode~sp.Rational(2,9)~ instead.
%
 \item\label{pi} The constant $\pi\simeq 3.1415926536$ should be entered as
  \pcode~sp.pi~.  If you just enter \pcode~pi~ then that will just 
  be treated as the name of a variable with no special value or
  properties. 
%
 \item\label{exp} The constant $e\simeq 2.718281828$ must be entered as 
  \pcode~sp.E~ (with a capital E) or \pcode~sp.exp(1)~.  
%
 \item\label{imaginary} The square root of minus one can be
  entered as \pcode~sp.I~.
%
 \item\label{infinity} In many contexts $\infty$ can be entered as
 \pcode~sp.oo~. 
\end{secenum}

\section{Solving}

\begin{secenum}
%
 \item\label{solve} To solve the equation $x^2+6=5x$ (for example), enter
  \pcode~sp.solve(sp.Eq(x**2+6,5*x))~.  If you want to solve for an
  expression being equal to zero then you do not need \pcode~sp.Eq()~,
  you can just enter \pcode~sp.solve(x**2+6-5*x)~ for example.  
  Similarly, if you have already defined a function $g$, you can enter
  \pcode~solve(g(t),t)~ to solve the equation $g(t)=0$. It
  \textbf{does not work} to enter \pcode~a == b~ instead of
  \pcode~sp.Eq(a, b)~.  
%
  \item\label{roots}
   The above syntax gives the answer in the form
   \pcode+[2,3]+.  If you prefer to have the answer in
   the form \pcode+[{x:2},{x:3}]+ (which can be more convenient if you
   want to substitute the result into some other expression), you
   should instead enter \pcode~sp.solve({x**2-5*x+6},dict=True)~.
%
  \item\label{realroots}
   The above syntax will give usually complex roots as well as real roots.  
   For example, if we declare \pcode+x = sp.symbols('x')+ and then enter 
   \pcode~sp.solve(x**3+x,x)~ we get $[0,-i,i]$.  If we only want real 
   roots we can either use \pcode~sp.realroots(x**3+x,x)~ or we can
   declare $x$ using \pcode+x = sp.symbols('x',real=True)+ and then
   just use \pcode+sp.solve()+.
%
 \item\label{multisolve} To solve several equations in several variables, put the
  equations in one set of square brackets, and the variables in another
  set.  For example, to solve $2x+3y=5$ and $3x+4y=7$ for $x$ and $y$,
  enter 
\begin{pcodeblock}
 sp.solve([2*x+3*y-5,3*x+4*y-7],[x,y],dict=True)
\end{pcodeblock}
%
 \item\label{nosols} If there are no solutions, \pcode~sp.solve()~ will
 return the empty list \pcode~[]~. 
%
 \item\label{allsolutions} If the problem involves trigonometric
  functions, then there will typically be infinitely many
  solutions, but \pcode~sp.solve()~ will only list a few of them.  
  To get all solutions, you can try using \pcode~sp.solveset()~
  instead.  For example, we can solve the equation $\cos(x)=1$ by entering
  \pcode+sp.solveset(sp.cos(x)-1,x)+.  The result will be printed as
  $\{2n\pi\;\mid\;n\in\mathbb{Z}\}$, indicating that $\cos(x)=1$ if
  and only if $x$ is an integer multiple of $2\pi$.  There is quite a
  long story about how this result is represented internally in
  Python, which you can read at 
  \url{https://docs.sympy.org/latest/modules/sets.html}.
%
 \item\label{fsolve} To get approximate numerical solutions, you can 
  use \pcode~sp.nsolve()~ instead of \pcode~sp.solve()~.  For example,
  you can enter \pcode~sp.nsolve(sp.cos(x)-sp.cos(x+1),x,42)~ to find
  a number $x$ with $\cos(x)=\cos(x+1)$, starting the search at
  $x=42$.  When using \pcode~sp.nsolve()~, you always need to specify
  the starting point.  If you need to solve a large number of
  equations numerically then it will usually be best to use the
  \pcode~scipy~ or \pcode~numpy~ packages; you should only use
  \pcode~sp.nsolve()~ for occasional numerical solutions when you are
  mostly working with symbolic equations.
% 
 \item\label{identity} Sometimes we need to solve problems like
  this: find $a$ and $b$ such that $a(x-1)^2+b(x+1)^2=x$
  \emph{for all x}.  For this we can enter 
  \begin{pcodeblock}
   sp.solve_undetermined_coeffs(a*(x-1)**2+b*(x+1)**2-x,[a,b],x)
  \end{pcodeblock}
  Similarly, to find $u$ such that $\sin(t+u)=\cos(t)$ for all
  $t$, we could try 
  \begin{pcodeblock}
   sp.solve_undetermined_coeffs(sp.Eq(sp.sin(t+u),sp.cos(t)),[u],t)
  \end{pcodeblock}
  This does not actually work, but we can make it work by giving the
  solver a hint as follows:
  \begin{pcodeblock}
   sp.solve_undetermined_coeffs(sp.Eq(sp.expand_trig(sp.sin(t+u)),sp.cos(t)),[u],t)
  \end{pcodeblock}
\end{secenum}

\section{Standard functions}

\begin{secenum}
%
 \item\label{standard} Most functions can be entered using their
  usual names prefixed with \pcode~sp.~, for example
 \pcode~sp.sin(x)~ or \pcode~sp.ln(y)~.  If you do not want to type
 the prefix, you can change your initial import statement to include 
 \begin{pcodeblock}
  from sympy import sin, cos, tan, exp, ln
 \end{pcodeblock}
 (for example) then you can just enter \pcode~sin(x)~ instead of
 \pcode~sp.sin(x)~. 
%
 \item\label{funcbrackets}
  However, you must always use brackets: enter
  \pcode~sp.tan(x)~ instead of \pcode~tan x~, and \pcode~sp.sin(2*x)~
  instead of \pcode~sin 2x~. 
%
 \item\label{abs} The absolute value of $x$, traditionally written as
  $|x|$, must be entered as \pcode~sp.abs(x)~. 
%
 \item\label{log} Both \pcode~sp.ln(x)~ and \pcode~sp.log(x)~ refer to the natural
  logarithm (to base $e$) of $x$. 
%
 \item\label{logbase}
  If you want to use the logarithm to base $10$ then you
  should enter \pcode~sp.log(x, 10)~.  Similarly, enter
  \pcode~sp.log(x, 2)~ for $\log_2(x)$ (the logarithm to base $2$) and so on. 
%
 \item\label{expx} The function $e^x$ should be entered as \pcode~sp.exp(x)~. 
%
 \item\label{hyp} The hyperbolic function $\sinh(x)=(e^x-e^{-x})/2$ can be
  entered as \pcode~sp.sinh(x)~, and similarly for the functions
  $\cosh(x)=(e^x+e^{-x})/2$ and $\tanh(x)=\sinh(x)/\cosh(x)$.
%
 \item\label{sinsquared} The traditional notation $\sin^2(x)$
  refers to the square of $\sin(x)$, which could also be
  written $(\sin(x))^2$ or $\sin(x)^2$.  In Python this must 
  be entered as \pcode~sp.sin(x)**2~.  Similar comments apply to
  $\tan^2(x)$ and so on.  Note that $\sin(x^2)$ means something
  different again: it is the sin of the square of $x$, not the square
  of the sin. 
%
 \item\label{cosec} In traditional notation $1/\sin(x)$ is sometimes written as
  $\text{csc}(x)$ or $\text{cosec}(x)$.  In Python \pcode~sp.csc(x)~ will
  work, but \pcode~sp.cosec(x)~ will not.  Of course, you can also just
  enter \pcode~1/sp.sin(x)~. 
%
 \item\label{arctan} In traditional notation $\tan^{-1}(x)$ or $\arctan(x)$
  or $\text{atan}(x)$ refers to the angle $\theta$ such that
  $\tan(\theta)=x$.  This is completely different from the number
  $\tan(x)^{-1}=1/\tan(x)$.  In Python you should enter
  \pcode~sp.atan(x)~ (\textbf{not} \pcode~tan^(-1)(x)~ or anything
  like that).  Similar comments apply 
  to $\text{asin}(x)$, $\text{atanh}(x)$ and so on. 
%
 \item\label{possqrt} When $x\geq 0$, the notation \pcode~sp.sqrt(x)~ means, by
  definition, the positive square root of $x$.  (There is more to say
  if $x$ is negative or complex, but we pass over that here.)  Thus,
  for example, \pcode~sp.sqrt(9)~ is definitely $3$ and not $-3$.  
\end{secenum}

\section{Defining your own functions}\label{sec-funcdef}

\begin{secenum}
%
 \item\label{arrow} If you want to define a simple function like 
  $f(x)=x^2$, enter \pcode~f = lambda x: x^2~.  It
  \textbf{does not} work to enter \pcode~f(x)=x^2;~ instead.
%
 \item\label{seriesdef}
  Similarly, to define the sequence $a(n)=(-1)^n/n$, enter
  \pcode~a = lambda n: (-1)**n/n~, not \\
  \pcode~a(n)=(-1)**n/n~. 
%
 \item\label{badderiv} It \textbf{does not} work to use notation like
   \pcode~f'(x)~ for the derivative of $f(x)$; see
   Section~\ref{sec-diff} instead. 
%
 \item\label{multifunc}
  You can define a function of several variables in a similar
  way.  For example, to define $g(a,u,v)=u^a+v^a$, enter
  \pcode~g = lambda a,u,v: u**a+v**a~. 
%
 \item\label{longfunc}
  For longer and more complicated functions you may prefer to use
  the more standard Python function definition syntax, like
  \begin{pcodeblock}
   def f(x, y, z):
       return (x+y)**8 + (x-y)**8 + (x+z)**8 + (x-z)**8
  \end{pcodeblock}
\end{secenum}

\section{Plotting}

In this course we mostly use the \pcode+sympy+ library, but that does
not, by itself, support plotting.  We instead use the \pcode+numpy+
and \pcode+plotly+ libraries to make plots.  (There is also another
popular library called \pcode+matplotlib+ which could be used instead
of \pcode+plotly+, but on balance we have found \pcode+plotly+ to be
preferable.)  More details of everything discussed here can be found
at \url{https://plotly.com/python/}.

\begin{secenum}
%
 \item\label{basicplot} To plot the graph $y=x^3-x$ from $x=-2$ to $x=2$ (for
  example), we would enter
  \begin{pcodeblock}
   x = np.linspace(-2, 2, 201)
   y = x**3 - x
   fig = go.Figure()
   fig.add_trace(go.Scatter(x=x, y=y))
   fig.show()
  \end{pcodeblock}
  The first line sets $x$ to be an array
  $[x_0,x_1,x_2,\dotsc,x_{200}]=[-2,-1.98,-1.96,\dotsc,2]$ consisting
  of $201$ equally spaced points between $-2$ and $2$ (inclusive).
  The second line sets $y$ to be the array
  $[y_0,\dotsc,y_{200}]=[x_0^3-x_0,\dotsc,x_{200}^3-x_{200}]$.  The
  third line creates an object called \pcode+fig+ representing a blank
  picture, to which we can add curves, shapes text and so on.  In this
  case we just add a single curve passing through the points with
  coordinates $(x_i,y_i)$; this is done by the line
  \pcode+fig.add_trace(go.Scatter(x=x, y=y))+.  Finally, the line
  \pcode+fig.show()+ to show the picture that we have constructed.
  This is not always necessary, if the last line in a cell is related
  to plotting then Python will often decide for itself that the
  resulting plot should be shown, but it does not hurt to make it
  explicit.
% 
 \item\label{legend} To label the curve in the above example, we could
  change the relevant line to 
  \begin{pcodeblock}
   fig.add_trace(go.Scatter(x=x, y=y, name='y=x^3-x', showlegend=True))
  \end{pcodeblock}
  This will display the label as \pcode+y=x^3-x+.  If we prefer to have 
  nicely formatted mathematics like $y=x^3-x$ then we can change the
  above to say \pcode+name=r'$y=x^3-x$'+.  However, this will only work
  correctly if we have done \pcode+from header import *+ as explained
  in Section~\ref{sec-start}.
% 
 \item\label{linestyle} Various features of the plotted curve can be changed
  by adding extra arguments inside \pcode+go.Scatter()+.  For example, 
  if we enter
  \begin{pcodeblock}
   fig.add_trace(go.Scatter(x=x, y=y, line_color='red', line_dash='dot', line_width=5))
  \end{pcodeblock}
  then the curve will be shown as a thick red dotted line.  Colours
  can be entered as names like \pcode~orange~
 or as hex codes (like
  \pcode+#40e0d0+ for turquoise; an internet search will find other
  examples).  Note that American spelling (color not colour) is
  required. 
%
 \item\label{layout} Various features of the overall picture can be changed 
  using \pcode+fig.update_layout()+ (before \pcode+fig.show()+).  For
  example, we can insert
  \begin{pcodeblock}
   fig.update_layout(height=400, width=800, xaxis_range=[-3,3], yaxis_range=[0,10])
  \end{pcodeblock}
  to ensure that the plot is $800\times 400$ pixels and shows $x$
  values from $-3$ to $3$ and $y$ values from $0$ to $10$.
%
 \item\label{scaling} To make the vertical scale the same as the horizontal
  scale, add the options\\
  \pcode~yaxis_scaleanchor="x",yaxis_scaleratio=1~ 
  in \pcode+fig.update_layout()+.  Alternatively, call
  \pcode+equal_aspect(fig)+ before \pcode+fig.show()+.
  (This is not a standard part of plotly and
  relies on having called \pcode+from header import *+ at the top of
  the notebook.)
%
 \item\label{twoplots} To plot more than one graph in the same picture, we
  can just call \pcode+fig.add_trace()+ several times.  For example, we
  can plot $y=x^3-x$ and $y=x^3+x$ in the same picture as follows:
  \begin{pcodeblock}
   x = np.linspace(-2, 2, 201)
   y0 = x**3 - x
   y1 = x**3 + x
   fig = go.Figure()
   fig.add_trace(go.Scatter(x=x, y=y0))
   fig.add_trace(go.Scatter(x=x, y=y1))
   fig.show()
  \end{pcodeblock}
  If we do not specify the colours, they will be chosen automatically
  so that the two curves have different colours.  Similarly,
  to plot $\cos(t)$ and $\sin(t)$ 
  together from $t=-3\pi$ to $t=3\pi$, enter
  \begin{pcodeblock}
   t = np.linspace(-3*np.pi, 3*np.pi, 201)
   fig = go.Figure()
   fig.add_trace(go.Scatter(x=t, y=np.cos(t)))
   fig.add_trace(go.Scatter(x=t, y=np.sin(t)))
   fig.show()
  \end{pcodeblock}
%
 \item\label{piticks} In the above example, there are tick marks on the
  $x$-axis at $x=-8,\;x=-7,\;\dotsc,x=8$.  It would be more natural to
  have the tick marks at $x=-3\pi,\;x=-2\pi,\;\dotsc,x=3\pi$.  We can
  achieve this by calling \pcode+use_pi_ticks(fig,-3,3)+ before
  \pcode+fig.show()+.  (This is not a standard part of plotly and
  relies on having called \pcode+from header import *+ at the top of
  the notebook.)  For ticks with a spacing of $\pi/2$ instead of $\pi$,
  we can call \pcode+us_pi_ticks(fig,-3,3,2)+.
 \item\label{hideaxes} To make the $x$ and $y$ axes invisible (together with
  the associated tick marks and labels) add the options
  \pcode+xaxis_visible=False, yaxis_visible=False+ in
  \pcode+fig.update_layout()+. 
 \item\label{discont} If we plot a function with discontinuities, like
  $y=1/x$, then by default we will get spurious vertical lines at the
  discontinuity points.  This can be avoided as follows:
  \begin{pcodeblock}
   x = np.linspace(-2,2,400)
   y = 1/x
   jumps = np.where(np.abs(np.diff(y)) >= 0.5)[0]
   x = np.insert(x, jumps+1, np.nan)
   y = np.insert(y, jumps+1, np.nan)
   fig = go.Figure()
   fig.add_trace(go.Scatter(x=x, y=y, mode='lines', 
                            name=r'$y = \{x\}$', showlegend=True))
   fig.update_layout(height=400, width=800, yaxis_range=[-8,8], showlegend=True)   
  \end{pcodeblock}
  The third line finds the values of $x$ where $y$ jumps by more than
  $0.5$, indicating a discontinuity.  (The value $0.5$ might need to
  be adjusted to give the right outcome in other examples.)  The next
  two lines insert the value \pcode~np.nan~ (indicating ``not a
  number'') at the corresponding places in the arrays $x$ and $y$.
  The plotting system takes this as an indication that the curve
  should be broken without spurious vertical lines.
 \item\label{numpoints} To improve the accuracy of a graph, just increase
  the third argument in \pcode+np.linspace+.  For example, the code
  \begin{pcodeblock}
   x = np.linspace(-1, 1, 50)
   fig = go.Figure()
   fig.add_trace(go.Scatter(x=x, y=x * np.sin(np.pi / x)))
  \end{pcodeblock}
  gives a very jagged picture, but it becomes much better if we change 
  \pcode+np.linspace(-1,1,50)+ to \pcode+np.linspace(-1,1,1000)+.
%
 \item\label{parametric} To plot the curve given parametrically by
  $x=1/(1+t^2)$ and $y=\sin(t)$ from $t=-5$ to $t=5$ (for example), enter
  \begin{pcodeblock}
   t = np.linspace(-5, 5, 400)
   x = 1/(1 + t ** 2)
   y = np.sin(t)
   fig = go.Figure()
   fig.add_trace(go.Scatter(x=x, y=y))
  \end{pcodeblock}
  Similarly, to plot the curve $(x,y)=(\sin(2t),\cos(3t))$
  for $t=0$ to $2\pi$, enter
  \begin{pcodeblock}
   t = np.linspace(0, 2*np.pi, 400)
   fig = go.Figure()
   fig.add_trace(go.Scatter(x=np.sin(2*t), y=np.cos(3*t)))
  \end{pcodeblock}
 \item\label{implicitplot} Consider a curve given implicitly, say
  by the equation $x+y+x^2y^2=1$ with $-3\leq x\leq 3$ and
  $-4\leq y\leq 4$.  This can be plotted as follows:
  \begin{pcodeblock}
   x0 = np.linspace(-3,3,100)
   y0 = np.linspace(-4,4,100)
   x, y = np.meshgrid(x0, y0)
   z = x + y + x**2 * y**2
   fig = go.Figure()
   fig.add_trace(go.Contour(x=x0, y=y0, z=z, 
                            contours=dict(type='constraint', value=1), 
                            line_color='red'))
   fig.update_layout(height=600, width=600, xaxis_range=[-3,3], yaxis_range=[-4,4])
  \end{pcodeblock}
  The first line sets \pcode+x0+ to ba a one-dimensional array of
  $x$-values, and similarly for the second line.  The third line sets
  $x$ to be a two-dimensional array of $x$-values, and $y$ to be a
  two-dimensional array of $y$-values.  The fourth line uses these to
  compute a two-dimensional array $z$ of values of the function
  $x+y+x^2y^2$, then the sixth line draws the curve where $z=1$.
  Similarly, you can plot the circle $x^2+y^2=4$ like this:
  \begin{pcodeblock}
   x0 = np.linspace(-2,2,100)
   y0 = np.linspace(-2,2,100)
   x, y = np.meshgrid(x0, y0)
   z = x**2 + y**2 - 1
   fig = go.Figure()
   fig.add_trace(go.Contour(x=x0, y=y0, z=z, 
                            contours=dict(type='constraint', value=0), 
                            line_color='red'))
   fig.update_layout(height=600, width=600, xaxis_range=[-2,2], yaxis_range=[-2,2])
  \end{pcodeblock}
%
 \item\label{saveplot} We might want to produce a basic picture, and then
  later generate some further plots which combine the basic one with
  some extra elements.  For this we can make the basic picture by
  calling \pcode+basic_fig=go.Figure()+ and then
  \pcode+basic_fig.add_trace(...)+.  Then we can call 
  \pcode+new_fig=go.Figure(basic_fig)+ to make a new figure which is a
  copy of the basic one, and \pcode+new_fig.add_trace(...)+ to add
  extra elements to the new figure, and \pcode+new_fig.show()+ to
  display the result.
%
 \item\label{display}
  Alternatively, if we have already created figures \pcode+fig0+ and
  \pcode+fig1+, we can make a new figure that combines them by
  entering \pcode+fig = go.Figure(data = fig0.data + fig1.data)+.
%
 \item\label{listplot}
  By default, \pcode+go.Scatter(x=u, y=v)+ will produce a connected
  curve passing through the points $(u_i,v_i)$.  Sometimes we might
  instead just want to show the points $(u_i,v_i)$ without connecting
  points.  For this, we need to add the option \pcode+mode='markers'+
  to \pcode+go.Scatter()+.  For example, we can display the points 
  $(11,50),(22,40),(33,30),(44,20),(55,10)$ as follows:
  \begin{pcodeblock}
   fig = go.Figure()
   fig.add_trace(go.Scatter(x=[11,22,33,44,55],y=[50,40,30,20,10], mode='markers'))
  \end{pcodeblock}
  Often we just want to display the $y$-values, and the $x$-values are
  not important, we just take them to be $0,1,\dotsc,n-1$.  In that
  context, we do not actually need to specify the $x$-values, we can
  just call something like 
  \pcode+fig.add_trace(go.Scatter(y=[50,40,30,20,10], mode='markers'))+
  and the $x$-values will be inserted automatically.
\end{secenum}

\section{Differentiation}\label{sec-diff}


\begin{secenum}
%
 \item\label{diff} To differentiate $x+\sin(x)$ with respect to $x$ (for
  example), enter \pcode~diff(x+sin(x),x)~.  This is
  equivalent to $\frac{d}{dx}(x+\sin(x))$ in traditional
  notation.  Similarly, to differentiate $1/(1-t^2)$ with
  respect to $t$ we enter \pcode~diff(1/(1-t^2),t)~.  
%
 \item\label{multidiff} To differentiate $y$ three times with respect to $x$
  (for example), enter \pcode~diff(y,x,3)~.  This is equivalent to
  $d^3y/dx^3$ in traditional notation.  Similarly, to find the second
  derivative of $\sin(\theta)$ with respect to $\theta$ we
  enter \pcode~diff(sin(theta),theta,2)~. 
%
 \item\label{implicitdiff} If $x$ and $y$ are related implicitly
  by an equation (say $x+y+x^2y^2=1$) then you can find
  $dy/dx$ in terms of $x$ and $y$ by implicit
  differentiation, like this: 
\begin{pcodeblock}
 sp.idiff(x+y+x^2*y^2-1,y,x)
\end{pcodeblock}
%
 \item\label{inertdiff} Sometimes we just want to mention the derivative
  without actually working it out.  We can do this by
  adding the option \pcode+evaluate=False+.  For example, 
  entering \\ \pcode+u = sp.diff(x**2, x, evaluate=False)+ gives an
  object $u$ which is displayed as $\frac{d}{dx}x^2$.  If you then
  enter \pcode+u.doit()+ then sympy will evaluate the derivative and
  return $2x$.  Similarly, \pcode~v = sp.diff(t**10,t,2)~ prints as
  $\frac{d^2}{dt^2}t^{10}$, whereas \pcode+v.doit()+ or
  \pcode~sp.diff(t**10,t,2)~  gives $90 t^8$. 
%
 \item\label{taylor} To find the Taylor series of a function $f(x)$ near a
  point $x=a$ to order $n$, enter
  \pcode~sp.series(f(x),x,a,n)~.  Note that ``order $n$'' means
  that the highest term allowed is a multiple of $x^{n-1}$,
  and all terms $x^n$ and higher are discarded.  For
  example, \pcode~sp.series(1/x,x,2,7)~ gives the Taylor series
  for $1/x$ near $x=2$ including terms up to $(x-2)^6$. 
%
  \item\label{taylorconv} The command \pcode~p = sp.series(...)~ command
  gives an answer including a term like $O(x^7)$, to indicate that
  there are infinitely many more terms, starting with a multiple of
  $x^{7}$.  We cannot do very much with this answer until we have
  converted it to an ordinary polynomial (without the $O(x^7)$ term on
  the end).  We can do this by entering \pcode+p = p.removeO()+.
\end{secenum}

\section{Integration}\label{sec-int}


\begin{secenum}
%
 \item\label{int} To calculate an indefinite integral like
  $\int\tan(x)\,dx$, enter \pcode~sp.integrate(sp.tan(x),x)~. 
%
 \item\label{intconst} Sympy's answer will never include an arbitrary
  constant '$+c$'.  In some contexts (particularly, solution
  of differential equations) the constant makes a
  difference and must be included.  In other contexts it
  is irrelevant; sympy cannot help you to decide about
  this. 
%
 \item\label{defint} To calculate a definite integral like
  $\int_1^5t^3\,dt$, enter \pcode~sp.integrate(t**3,(t,1,5))~. 
  Similarly, to calculate $\int_0^2\sqrt{1+x}\,dx$, enter
  \pcode~sp.integrate(sp.sqrt(1+x),x=0..2)~. 
%
 \item\label{improper} It is also possible to have infinite limits.  For
  example, $\int_{-\infty}^\infty (1+x^4)^{-1}dx$ can be 
  entered as 
  \begin{pcodeblock}
   sp.integrate(1/(1+x**4),(x,-sp.oo,sp.oo))
  \end{pcodeblock}
%
 \item\label{inertint} As with differentiation, it is sometimes useful to
  mention an integral without evaluating it.  You can do
  this by writing \pcode~sp.Integral()~ instead of
  \pcode~sp.integrate()~.  For example, if you enter
  \pcode~u = sp.Integral(x^2,x)~ then you will get an object that is
  displayed as $\int x^2\,dx$; if you then enter
  \pcode~u.doit()~, you will get $x^3/3$. 
\end{secenum}

\section{Sequences and sums}

\begin{secenum}
%
 \item\label{seq} To generate a finite sequence of terms, we can use
 standard Python list comprehension.  For example, the sequence
 $1/2,1/4,1/6,1/8,1/10$ is the sequence of terms $1/(2k)$ for
 $k=1,2,3,4,5$; it can be generated by the command 
 \begin{pcodeblock}
  list(sp.Rational(1,2*k) for k in range(1,6))
 \end{pcodeblock}
 Remember here that \pcode+range(1,6)+ refers to the range from $1$ to
 $6$ including $1$ but not including $6$, so it gives $1,2,3,4,5$ as
 required. 
%
 \item\label{sum} For a finite list $u$ as above, we can just enter
 \pcode+sum(u)+ to get the sum of all the entries in the list.  For
 example, \pcode+u=list(10**k for k in range(5))+ gives
 $u=[1,10,100,1000,10000]$ and then \pcode+sum(u)+ gives $11111$.
 Moreover, \pcode+sum(10**k for k in range(5))+ will also work for
 this; it is not necessary to make a list as an intermediate step.
%
 \item\label{symbolicsum} However, we might want to get a general formula
 for a sum like $\sum_{k=0}^n3^{-k}$ in terms of $n$, and this cannot 
 be done by the above method.  Instead, we need to declare $n$ and $k$
 to be integer symbols and use \pcode+sp.Sum()+, like this:
 \begin{pcodeblock}
  k, n = sp.symbols('k n', integer=True)
  u0 = sp.Sum(3**(-k),(k,0,n)).doit()
 \end{pcodeblock}
 Note that the naming and behaviour is a bit inconsistent with the
 naming and behaviour for integrals.  The function \pcode+sp.Sum()+
 always returns an unevaluated sum which is displayed as something
 like $\sum_{k=0}^n3^{-k}$; it is always necessary to use
 \pcode+.doit()+ to find the actual value of the sum.  Also, the range
 \pcode+(k,0,n)+ means that $k$ runs from $0$ to $n$ inclusive, which
 is consistent with tradional mathematical notation, but not with the
 usual Python convention for ranges.
%
 \item\label{infinitesum} We can also handle some infinite sums.  For
 example, it is a well-known result that
 $\sum_{k=1}^\infty k^{-2}=\pi^2/6$, and we can do this in sympy by
 entering \pcode+sp.Sum(k**(-2),(k,1,sp.oo)).doit()+.
\end{secenum}

\section{Some common problems}\label{sec-problems}

\begin{secenum}
 \item\label{symbolsagain} All abstract symbols need to be introduced using somthing
    like \pcode+x = sp.symbols('x')+.  If you try to use $x$ without
    having declared it, you will get an error like
    \pcode+name 'x' is not defined+.  It can also happen that you give
    $x$ a value, perhaps by entering $x = 42$, and later you forget
    that you did that, and try to treat $x$ as an abstract symbol.  
    This might give you any of a range of different error messages, or
    it might just give you an answer that is incorrect or meaningless
    but with no error message.  You can check the status of $x$ by
    entering \pcode+type(x)+.  If $x$ has been declared as an abstract
    symbol, then \pcode+type(x)+ should give
    \pcode+sympy.core.symbol.Symbol+.  If $x$ has been set equal to
    $42$, then \pcode+type(x)+ should give \pcode+int+.  If $x$ has
    not been introduced in any way, then \pcode+type(x)+ will give an
    error.  If $x$ has been given a value that is no longer relevant,
    you can enter \pcode+del x+ to remove it.
%
 \item\label{rememberstar} Remember \pcode~*~'s for multiplication.  
  \begin{itemize}
   \item Enter $2(x+y)$ as \pcode~2*(x+y)~.  If you instead enter 
    \pcode~2(x+y)~ then you may get a cryptic error message saying 
    \pcode+'int' object is not callable+.
   \item Enter $xy$ as $x*y$.  If you instead enter \pcode+xy+ then
    you will usually get an error message
    \pcode+name 'xy' is not defined+.
   \item For the product of $a$ and $x+y$, enter \pcode~a*(x+y)~.
    To evaluate a function $f(t)$ at the point $t=x+y$, enter
    $f(x+y)$.  If $a$ is an object for which \pcode~a*(x+y)~ makes
    sense, then $a(x+y)$ will not make sense and will give an error.
    If $f$ is an object for which $f(x+y)$ makes sense, then $f*(x+y)$
    will not make sense and will give an error.
  \end{itemize}
%
 \item\label{fracbrac} The fraction $\frac{a+b}{c+d}$ must be entered as
  \pcode~(a+b)/(c+d)~ (not \pcode~a+b/c+d~), and
   $\frac{a}{(u+v)(x+y)}$ must be entered as
   \pcode~a/((u+v)*(x+y))~ (not \pcode~a/(u+v)*(x+y)~). 
%
 \item\label{expbrac} $a^{1/4}$ must be entered as
  \pcode~a ** sp.Rational(1,4)~, not \pcode+a ** (1/4)+ or
  \pcode~a ** 1/4~. 
%
 \item\label{shape} Always use round brackets (not square or curly
         ones) for algebraic grouping.  Square brackets are
         used for lists or for subscripts (e.g. $x_5$ is
         \pcode~x[5]~).  Curly brackets are used for sets. 
%
 \item\label{removedefs} Remember to remove definitions when you
  have finished with  them, by using \pcode+del+ or clicking the
  \pcode~Restart;~ button.  If you enter an expression like $(x+y)^5$
  and SymPy converts it to something completely unrelated like
  $(\sin(\theta)+e^{-u})^5$, the most likely reason is that
  you previously defined \pcode~x=sp.sin(theta);~ and
  \pcode~y=sp.exp(-u);~, and you forgot to remove these
  definitions.  You can enter 
  \pcode+x,y = sp.symbols('x y')+ to reset $x$ and $y$ to be 
  symbols with no assigned value.
%
 \item\label{badexp} Remember the constant $e=2.71828\dots$ must be
  entered as \pcode+sp.E+ or \pcode~sp.exp(1)~, and $e^x$ must be entered as
  \pcode~exp(x)~.  If you enter \pcode+e ** x+ instead of \pcode+sp.exp(x)+,
  you will most often get an error message
  \pcode+name 'e' is not defined+.  However, if you happen to have
  defined $e$, you might get some other behaviour.
%
 \item\label{badpiagain} Remember that the constant
  $\pi=3.14159\dots$ must be 
  entered as \pcode~sp.pi~, not \pcode~pi~.  If sympy is doing
  funny things like refusing to simplify $\sin(\pi)$ to $0$,
  it may be that you entered \pcode~pi~ instead of \pcode~sp.pi~
  by mistake.  You can cure this by entering \pcode~pi=sp.pi~.
\end{secenum}

\end{document}
